% =====================================================================
%  COVER LETTER — example for a FICTIONAL person (Alex Sample).
%  Compile: pdflatex cover-letter.tex   (run twice)
%
%  HOW TO TAILOR THIS FILE
%   1. \cvtitle must match the CV and the posting's exact job title.
%   2. Fill the recipient block; use the named hiring manager if the
%      posting gives one, otherwise "Dear Hiring Team,".
%   3. Three paragraphs, no more, in this order:
%        P1 — why this company and this role, drawn from the actual ad
%        P2 — professional match against the posting's real tasks,
%             closing with one honest gap-acknowledgment sentence
%        P3 — what you are like as a colleague, grounded in real facts
%   4. Target 300-380 words of body text. It must stay on one page.
% =====================================================================
\documentclass[12pt,a4paper]{article}
\usepackage{jobcv}

% ------------------------------------------------------------ identity
\cvname{Alex Sample, M.Sc.}
\cvtitle{Embedded Systems Engineer}
\cvlocation{Utrecht, NL}
\cvphone{+31 6 0000 0000}
\cvemail{alex.sample@example.com}
\cvsite{alexsample.example}{https://alexsample.example}
\cvlinkedin{linkedin.com/in/alex-sample}{https://www.linkedin.com/in/alex-sample/}
% \cvphoto{headshot.jpg}

\begin{document}
\pagestyle{empty}
\makecvheader

% ---------------------------------------------------------- recipient
\cldate{Utrecht, 10 October 2026}

\clrecipient{%
Company Name\\
Hiring Manager Name, Job Title\\
Street Address\\
Postcode City, Country}

% ------------------------------------------------------------ subject
\clsubject{Application: Embedded Systems Engineer}

\clsalutation{Dear Hiring Team,}

\begin{clbody}
Open with why this company and this role, using specifics from the posting itself rather than
generic praise. Name the product, the team, or the problem the advertisement actually describes,
and say what draws you to it. Two or three sentences is enough --- the reader is deciding whether
to keep reading, not grading your enthusiasm.

Match your experience against the tasks the posting actually lists, in the posting's own
vocabulary, and ground every claim in something real: a system you built, a number you measured,
a method you ran. Lead with the closest match rather than the most impressive one. Then close
this paragraph by naming the one requirement you do not yet meet, plainly and without apology
--- for example, \enquote{I have not worked with X before, but I have moved into unfamiliar
technical domains several times and I get productive quickly.} Naming the gap yourself is more
convincing than hoping it goes unnoticed.

Finish with what you are like to work with, using facts rather than adjectives: teams led,
people mentored, cross-functional collaborations. If you have a portfolio page, link it once
here, tied to something you already argued --- for example, the sensor platform is documented
at \weblink{https://alexsample.example/projects}{alexsample.example/projects}. Close by inviting
the conversation.
\end{clbody}

\clclosing{Kind regards,}{Alex Sample}

\end{document}
